\section{結論}
\label{sec:conclusion}

本稿は，Apriori-window の枠組みを維持しつつ，入力データのバスケット構造を
明示的に扱う Basket-aware Apriori-Window（Phase 1）を提案した．
提案法は，共起判定をトランザクション内からバスケット内へ置き換えることで，
偽共起の原因となるバスケット横断混入を抑制する．

合成データ実験では，(i) バスケット数増加時に従来法の偽共起が急増すること（A1），
(ii) カテゴリ構造の変化に応じて偽共起率が変動すること（A2），
(iii) 大規模条件で提案法の実行時間優位が拡大すること（A4）を確認した．
特に $10^6$ トランザクションでは，提案法 95.5 秒，従来法 143.6 秒であり，
従来法は 1.50 倍遅かった．

今後の課題としては，(1) 実データ（Stage B）での外的妥当性評価，
(2) 区間品質指標（Purity/Jaccard）を含む総合評価，
(3) 並列化・メモリ管理の最適化によるさらなる高速化が挙げられる．
